% Part: proof-theory
% Chapter: sequent-calculus
% Section: quantifiers

\documentclass[../../../include/open-logic-section]{subfiles}

\begin{document}

\olfileid{pt}{seq}{qua}

\olsection{برهان‌های منظم و جانشینی}

قواعد سورها به‌سبب «شرط‌های متغیر ویژه» که بر قواعد
\LeftR{\lexists} و~\RightR{\lforall} حاکم‌اند، پیچیدگی‌هایی پدید
می‌آورند. بیشتر این پیچیدگی‌ها را می‌توان با این تضمین رفع کرد که
!!{constant}~$a$ در این استنتاج‌ها دوباره به کار نرود. تعریف زیر را
در نظر می‌گیریم.

\begin{defn}
!!^a{proof}~$\pi$ در~\Log{G1c} یا~\Log{G3c} را \emph{منظم} می‌نامیم
هرگاه:
\begin{enumerate}
  \item هیچ متغیر ویژه‌ای چون~$a$ در بیش از یک استنتاج
  \LeftR{\lexists} یا~\RightR{\lforall} به‌عنوان متغیر ویژه به کار
  نرفته باشد؛ و
  \item اگر~$a$ متغیر ویژهٔ یک استنتاج~\LeftR{\lexists} یا
  \RightR{\lforall} باشد، در~$\pi$ فقط بالای همان استنتاج رخ دهد.
\end{enumerate}
\end{defn}

\begin{lem}\ollabel{lem:subst-regular} اگر~$\pi$ منظم باشد و به
$\Gamma \Sequent \Delta$ ختم شود، $c$~!!a{constant} باشد که در~$\pi$
متغیر ویژه نیست، و~$t$ جمله‌ای باشد که هیچ متغیر ویژهٔ~$\pi$ را در بر
ندارد، آنگاه~$\Subst{\pi}{t}{c}$، که با جایگزین‌کردن هر رخداد~$c$ با
$t$ در~$\pi$ حاصل می‌شود، !!{proof}ی منظم است. دنبالهٔ پایانی آن
$\Subst{\Gamma}{t}{c} \Sequent \Subst{\Delta}{t}{c}$ است؛ در اینجا
$\Subst{\Gamma}{t}{c}$ چندمجموعه‌ای است که با جایگزینی~$c$ با~$t$ در
هر رخداد فرمولیِ~$\Gamma$، با حفظ همهٔ رخدادها و چندگانگی‌هایشان، به دست
می‌آید.
\end{lem}

\begin{proof}
  بر~$\pheight{\pi}$ استقرا می‌کنیم. اگر~$\pheight{\pi}=0$، $\pi$
  فقط از یک اصل تشکیل شده است. در این صورت~$\Subst{\pi}{t}{c}$ نیز اصل
  است، زیرا هر~!!{formula}~$!A(c)$ در آن به~$!A(t)$ تبدیل می‌شود.

  اگر~$\pheight{\pi}>0$ باشد، بر حسب آخرین استنتاج~$\pi$ حالت‌بندی
  می‌کنیم. حالت‌های قواعد ساختاری (در~\Log{G1c}) و قواعد منطقیِ رابط‌های
  گزاره‌ای آسان‌اند و به‌عنوان تمرین واگذار می‌شوند. حالت‌هایی را بررسی
  می‌کنیم که~$\pi$ به استنتاجی سورمند ختم می‌شود.
  \begin{enumerate}
    \item اگر آخرین استنتاج~\RightR{\lforall} باشد، $\pi$ به صورت
    زیر است:
    \[
    \AxiomC{}
    \RightLabel{$\pi'$}
    \Deduce$\Gamma(c) \fCenter \Delta'(c), !A(a,c)$
    \RightLabel{\RightR{\lforall}}
    \UnaryInf$\Gamma(c) \fCenter \Delta'(c), \lforall[x][!A(x,c)]$
    \DisplayProof
    \]
    بنا بر فرض استقرا، $\Subst{\pi'}{t}{c}$~!!{proof}ی منظم از
    $\Gamma(t) \fCenter \Delta'(t), !A(a,t)$ است. چون~$a$ متغیر ویژهٔ
    $\pi$ است، در~$t$ رخ نمی‌دهد. پس آخرین استنتاج در
    $\Subst{\pi}{t}{c}$،
    \[
    \AxiomC{}
    \RightLabel{$\Subst{\pi'}{t}{c}$}
    \Deduce$\Gamma(t) \fCenter \Delta'(t), !A(a,t)$
    \RightLabel{\RightR{\lforall}}
    \UnaryInf$\Gamma(t) \fCenter \Delta'(t), \lforall[x][!A(x,t)]$
    \DisplayProof
    \]
    درست است: نتیجه شامل~$a$ نیست.
    \item اگر آخرین استنتاج~\LeftR{\lforall} در~\Log{G3c} باشد،
    $\pi$ به صورت
    \[
    \AxiomC{}
    \RightLabel{$\pi'$}
    \Deduce$!A(s(c),c), \lforall[x][!A(x,c)], \Gamma(c) \fCenter \Delta'(c)$
    \RightLabel{\LeftR{\lforall}}
    \UnaryInf$\lforall[x][!A(x,c)], \Gamma(c) \fCenter \Delta'(c)$
    \DisplayProof
    \]
    است. بنا بر فرض استقرا، $\Subst{\pi'}{t}{c}$~!!{proof}ی منظم از
    $!A(s(t),t), \lforall[x][!A(x,t)], \Gamma(t) \fCenter
    \Delta'(t)$ است. آخرین استنتاج در~$\Subst{\pi}{t}{c}$،
    \[
    \AxiomC{}
    \RightLabel{$\Subst{\pi'}{t}{c}$}
    \Deduce$!A(s(t),t), \lforall[x][!A(x,t)], \Gamma(t) \fCenter \Delta'(t)$
    \RightLabel{\LeftR{\lforall}}
    \UnaryInf$\lforall[x][!A(x,t)], \Gamma(t) \fCenter \Delta'(t)$
    \DisplayProof
    \]
    درست است. حالت~\LeftR{\lforall} در~\Log{G1c} مشابه، اما اندکی
    ساده‌تر است.
    \item اگر آخرین استنتاج~\RightR{\lexists} یا~\LeftR{\lexists}
    باشد، بررسی آن تمرین است.
  \end{enumerate}
  در هر حالت، دنباله‌ای به پایان یک~!!{proof} منظم (یا در قواعد
  دومقدمه‌ای به~!!{proof}s منظم) افزوده‌ایم.
  در هر حالت، استنتاج‌های متغیر ویژه تغییر نمی‌کنند. چون~$c$ متغیر ویژه
  نبود و~$t$ هیچ متغیر ویژهٔ~$\pi$ را در بر ندارد، جانشینی هیچ متغیر
  ویژه‌ای را پایین‌تر از استنتاج ویژهٔ آن وارد نمی‌کند؛ فرض‌های استقرا نیز
  منظمی را حفظ می‌کنند. بنابراین~$\Subst{\pi}{t}{c}$ منظم است.
\end{proof}

\begin{prop}
اگر~$\pi$~!!a{proof} در~\Log{G1c} یا~\Log{G3c} باشد، !!{proof}ی منظم
چون~$\pi'$ با همان ساختار (و به‌ویژه همان~!!{height}) وجود دارد که فقط
در جایگزینی متغیرهای ویژه با~$\pi$ تفاوت دارد.
\end{prop}

\begin{proof}
فقط برای این برهان، استنتاج متغیر ویژه‌ای از~$\pi$ را \emph{پاک}
می‌نامیم اگر متغیر ویژهٔ آن، متغیر ویژهٔ هیچ استنتاج دیگری نباشد و در
$\pi$ جز در~!!{proof} بی‌واسطهٔ آن استنتاج رخ ندهد؛ در غیر این صورت آن
را \emph{ناپاک} می‌نامیم. فرض کنید~$n$ شمار استنتاج‌های ناپاک متغیر ویژه
در~$\pi$ باشد. با استقرا بر~$n$ نشان می‌دهیم که~$\pi$ را می‌توان به
!!{proof} منظم مطلوب~$\pi'$ تبدیل کرد. اگر~$n=0$ باشد، همهٔ استنتاج‌های
متغیر ویژه در~$\pi$ پاک‌اند؛ پس~$\pi$ منظم است و چیزی برای اثبات باقی
نمی‌ماند. در غیر این صورت، بالاترین استنتاج ناپاک متغیر ویژه را برمی‌گزینیم؛
فرض کنیم استنتاجی با~\RightR{\lforall} باشد. زیرـ!!{proof}~$\pi_1$ که
به آن ختم می‌شود، به صورت
    \[
    \AxiomC{}
    \RightLabel{$\pi_2$}
    \Deduce$\Gamma \fCenter \Delta, !A(c)$
    \RightLabel{\RightR{\lforall}}
    \UnaryInf$\Gamma \fCenter \Delta, \lforall[x][!A(x)]$
    \DisplayProof
    \]
است. چون~$c$ متغیر ویژه است، در~$\Gamma$، $\Delta$ یا
$\lforall[x][!A(x)]$ رخ نمی‌دهد. برهان~$\pi_2$ منظم است، زیرا هر
استنتاج متغیر ویژهٔ بالای استنتاج برگزیده پاک است. فرض کنید~$a$
!!a{constant} باشد که در~$\pi$ رخ نمی‌دهد. به موجب
\olref{lem:subst-regular}، $\Subst{\pi_2}{a}{c}$ برهانی منظم از
$\Gamma \fCenter \Delta, !A(a)$ است و می‌توان~\RightR{\lforall} را
بر آن اعمال کرد. اگر~$\pi_1$ را در~$\pi$ با برهان~$\pi_1'$،
    \[
    \AxiomC{}
    \RightLabel{$\Subst{\pi_2}{a}{c}$}
    \Deduce$\Gamma \fCenter \Delta, !A(a)$
    \RightLabel{\RightR{\lforall}}
    \UnaryInf$\Gamma \fCenter \Delta, \lforall[x][!A(x)]$
    \DisplayProof
    \]
جایگزین کنیم، متغیر ویژهٔ تازهٔ~$a$ را جایگزین~$c$ کرده‌ایم. برهان
حاصل اکیداً کمتر از~$n$ استنتاج ناپاک متغیر ویژه دارد، زیرا این تغییر نام
ممکن است بیش از یک استنتاج را پاک کند. اکنون نتیجه از فرض استقرا به دست
می‌آید. حالت‌های دیگر قواعد متغیر ویژه مشابه‌اند.
\end{proof}

از گزاره‌ای که همین اکنون اثبات شد صریحاً استفاده نخواهیم کرد، اما فرض
می‌کنیم همهٔ !!{proof}s مورد بررسی منظم‌اند؛ اگر چنین نباشند، همیشه
می‌توان متغیرهای ویژه را جایگزین کرد و آن‌ها را منظم ساخت.

\end{document}
